% Part: model-theory
% Chapter: basics
% Section: nonstandard-arithmetic

\documentclass[../../../include/open-logic-section]{subfiles}

\begin{document}

\olfileid{mod}{bas}{nsa}
\olsection{مدل‌های نااستاندارد حساب}

\begin{defn}
بگذارید $\Lang{L_N}$، !!{language} حساب باشد که از یک !!{constant}
~$\Obj{0}$، یک !!{predicate} دوموضعی $<$، یک !!{function} یک‌موضعی
$\prime$ و !!{function}های دوموضعی $+$ و $\times$ تشکیل شده است.
\begin{enumerate}
\item \emph{مدل استاندارد} حساب، !!{structure}~$\Struct{N}$ برای
  $\Lang{L_N}$ است که $\Nat = \{ 0, 1, 2, \dots\}$ دارد و
  $\Obj{0}$ را به‌صورت $0$، $<$ را به‌صورت رابطهٔ کوچک‌تربودن
  بر~$\Nat$، و $\prime$، $+$ و $\times$ را به‌ترتیب به‌صورت جانشین،
  جمع و ضرب بر $\Nat$ تفسیر می‌کند.
\item \emph{حساب صادق} نظریهٔ $\Theory{N}$ است.
\end{enumerate}
\end{defn}

هنگام کار در $\Lang{L_N}$، هر ترم به شکل
$\Obj{0}^{\prime\dots\prime}$ را که در آن تابع جانشین $n$ بار بر
$\Obj{0}$ اعمال شده است، با~$\num{n}$ کوتاه می‌کنیم.

\begin{defn}
یک !!{structure}~$\Struct{M}$ برای $\Lang{L_N}$ \emph{استاندارد} است
اگر و تنها اگر $\Struct{N} \iso \Struct{M}$.
\end{defn}

\begin{thm}\ollabel{thm:non-std}
مدل‌های !!{enumerable} نااستاندارد از حساب صادق وجود دارند.
\end{thm}

\begin{proof}
$\Lang{L_N}$ را با افزودن !!{constant} تازهٔ~$c$ گسترش دهید و نظریهٔ
\[
\Theory{N} \cup \Setabs{\num{n} < c}{n \in \Nat}
\]
را در نظر بگیرید. این نظریه به‌طور متناهی ارضاپذیر است؛ پس بنا بر
فشردگی مدلی مانند $\Struct{M}$ دارد که بنا بر قضیهٔ
لونهایم--اسکولم نزولی می‌توان آن را !!{enumerable} گرفت. اگر
$\Domain{M}$ دامنهٔ $\Struct{M}$ باشد، بگذارید $\Struct{M}$ نمادهای
غیرمنطقی $\Lang{L_N}$ را چنین تفسیر کند:
$\mathbf{z} = \Assign{\Obj{0}}{M} \in \Domain{M}$،
$\prec = \Assign{<}{M} \subseteq \Domain{M}^2$،
$* = \Assign{\prime}{M} \colon \Domain{M} \to \Domain{M}$، و
$\oplus = \Assign{+}{M}$ و
$\otimes = \Assign{\times}{M} \colon \Domain{M}^2 \to \Domain{M}$.
برای هر $x \in \Domain{M}$، عنصر $\Domain{M}$ که با اعمال~$*$ بر $x$
به دست می‌آید را $x^*$ می‌نویسیم.

اکنون اگر $h$ یک یکریختی میان $\Struct{N}$ و $\Struct{M}$ بود، عدد
$n \in \Nat$ای وجود داشت که $h(n) = \Assign{c}{M}$. پس بگذارید $s$
هر انتسابی در $\Struct{N}$ باشد که $s(x) = n$. آنگاه
$\Sat{N}{\eq[\num{n}][x]}[s]$؛ و بنا بر اثبات
\olref[iso]{thm:isom}،
$\Sat{M}{\eq[\num{n}][x]}[h \circ s]$ نیز. بنابراین
$\Assign{c}{M} = \mathbf{z}^{*\cdots *}$، که در آن $*$ به تعداد $n$
بار تکرار شده است. اما این ناممکن است، زیرا بنا بر فرض
$\Sat{M}{\num{n} < c}$ و $\prec$ ناتبازتابی است. پس $\Struct{M}$
نااستاندارد است.
\end{proof}

\begin{prob}
رابطهٔ $R$ بر مجموعهٔ $X$ \emph{خوش‌بنیاد} است اگر و تنها اگر هیچ
زنجیرهٔ نزولی نامتناهی در~$R$ نباشد؛ یعنی هیچ
$x_0$، $x_1$، $x_2$،~\dots در $X$ وجود نداشته باشند که
$\dots x_2Rx_1Rx_0$. با فرض سازگار‌بودن نظریهٔ مجموعه‌های
زرملو--فرانکل~$ZF$ نشان دهید مدل‌های ناخوش‌بنیاد از $ZF$ وجود دارند؛
یعنی مدل‌هایی مانند $\mathfrak{M}$ که
$\dots x_2 \in x_1 \in x_0$.
\end{prob}

چون مدل نااستاندارد $\Struct{M}$ از \olref{thm:non-std} به‌طور مقدماتی
با مدل استاندارد هم‌ارز است، می‌توان شماری از ویژگی‌های $\Struct{M}$
را به دست آورد. باقی این بخش به این کار اختصاص دارد و به ما امکان
می‌دهد مدل‌های نااستاندارد !!{enumerable} از $\Theory{N}$ را دقیقاً
توصیف کنیم.

\begin{enumerate}
\item هیچ عضوی از $\Domain{M}$ از خودش $\prec$-کوچک‌تر نیست: جملهٔ
  $\lforall[x][\lnot x < x]$ در $\Struct{N}$ و بنابراین در
  $\Struct{M}$ صادق است.
\item با استدلالی مشابه نتیجه می‌گیریم $\prec$ یک
  \emph{ترتیب خطی} بر $\Domain{M}$ است؛ یعنی رابطه‌ای تام، ناتبازتابی
  و متعدی بر $\Domain{M}$.
\item عنصر $\mathbf{z}$، $\prec$-کمینهٔ $\Domain{M}$ است.
\item هر عضو $\Domain{M}$ از جانشین $*$ خود $\prec$-کوچک‌تر است و
  $x^*$، $\prec$-کمینهٔ اعضای $\Domain{M}$ است که از $x$ بزرگ‌ترند.
\item $\Struct{M}$ یک پارهٔ آغازی، نسبت به $\prec$، دارد که با $\Nat$
  یکریخت است:
  $\mathbf{z}, \mathbf{z}^*, \mathbf{z}^{**}, \dots$.
  آن را \emph{بخش استاندارد} $\Domain{M}$ می‌نامیم. هر عضو دیگر
  $\Domain{M}$ \emph{نااستاندارد} است. باید اعضای نااستانداردی در
  $\Domain{M}$ وجود داشته باشند، وگرنه تابع $h$ از اثبات
  \olref{thm:non-std} یک یکریختی بود. از $n,m,\dots$ به‌عنوان
  !!{variable}هایی استفاده می‌کنیم که بر این بخش استاندارد
  $\Struct{M}$ تغییر می‌کنند.
\item هر عنصر نااستاندارد از هر عنصر استاندارد بزرگ‌تر است؛ زیرا برای
  هر $n \in \Nat$،
  \[
  \Sat{N}{\lforall[z][(\lnot(\eq[z][\Obj{0}] \lor
  \dots \lor \eq[z][\num{n}]) \lif \num{n} < z)]},
  \]
  و بنابراین اگر $z \in \Domain{M}$ با همهٔ عناصر استاندارد متفاوت
  باشد، باید از همهٔ آن‌ها \emph{بزرگ‌تر} باشد.
\item هر عضو $\Domain{M}$ جز $\mathbf{z}$ جانشین $*$ عنصری یکتا از
  $\Domain{M}$ است که با $^*x$ نشان داده می‌شود. اگر $x = y^*$،
  آنگاه اگر یکی از $x$ و $y$ استاندارد باشد هر دو استانداردند، و اگر
  یکی نااستاندارد باشد هر دو نااستانداردند.
\item رابطهٔ هم‌ارزی $\approx$ بر $\Domain{M}$ را چنین تعریف کنید:
  $x \approx y$ اگر و تنها اگر برای یک $n$ \emph{استاندارد}، یا
  $x \oplus n = y$ یا $y \oplus n =x$. به بیان دیگر، $x \approx y$
  اگر و تنها اگر $x$ و $y$ فاصله‌ای متناهی داشته باشند. اگر $n$ و
  $m$ استاندارد باشند، آنگاه $n \approx m$. \emph{بلوک} $x$ را کلاس
  هم‌ارزی $[x] = \Setabs{y}{x \approx y}$ تعریف کنید.
\item فرض کنید $x \prec y$ و $x \not\approx y$. چون
  $\Sat{N}{\lforall[x][\lforall[y][(x < y \lif (x' < y \lor x' =
      y))]]}$، یا $x^* \prec y$ یا $x^* = y$. حالت دوم ناممکن است،
  زیرا $x \approx y$ را نتیجه می‌دهد؛ پس $x^* \prec y$. به همین
  ترتیب، اگر $x \prec y$ و $x \not\approx y$، آنگاه
  $x \prec {^*y}$. بنابراین اگر $x \prec y$ و $x \not\approx y$،
  هر $w \approx x$ از هر $v \approx y$ به نسبت $\prec$ کوچک‌تر است.
  در نتیجه هر بلوک نااستاندارد $[x]$ زنجیره‌ای دوسویه و نامتناهی
  تشکیل می‌دهد:
  \[
  \cdots \prec  {^{**}x} \prec {^*}x \prec x \prec x^* \prec x^{**}
  \prec \cdots
  \]
  این زنجیره را $Z$-زنجیره می‌نامند، زیرا نوع ترتیبی اعداد صحیح را دارد.
\item ترتیب $\prec$ را می‌توان به بلوک‌ها بالا برد: اگر $x \prec y$ و
  $x \not\approx y$، آنگاه بلوک $x$ از بلوک $y$ کوچک‌تر است. بلوکی
  \emph{نااستاندارد} است اگر عنصری نااستاندارد در بر داشته باشد. بلوک
  استاندارد کوچک‌ترین بلوک است.
\item کوچک‌ترین بلوک نااستاندارد وجود ندارد: اگر $y$ نااستاندارد باشد،
  آنگاه یک $x \prec y$ وجود دارد که $x$ نیز نااستاندارد و
  $x \not\approx y$ است. اثبات: در مدل استاندارد $\Struct{N}$، هر عدد
  بر دو بخش‌پذیر است، شاید با باقیماندهٔ یک؛ یعنی
  $\Sat{N}{\lforall[y][\lexists[x][(y = x + x \lor y = x + x +
        \Obj{0}')]]}$. بنا بر هم‌ارزی مقدماتی، برای هر
  $y \in \Domain{M}$، یک $x \in \Domain{M}$ هست به‌طوری‌که یا
  $x \oplus x = y$ یا
  $x \oplus x \oplus \mathbf{z}^*= y$. اگر $x$ استاندارد بود، $y$ نیز
  استاندارد می‌بود؛ پس $x$ نااستاندارد است. افزون بر این، $x$ و $y$
  به بلوک‌های متفاوت تعلق دارند؛ یعنی $x \not\approx y$. برای دیدن
  این امر، فرض کنید در یک بلوک باشند؛ یعنی برای $n$ استانداردی
  $x \oplus n = y$. اگر $y = x \oplus x$، آنگاه
  $x \oplus n = x \oplus x$؛ پس بنا بر قانون حذف برای جمع، که در
  $\Struct{N}$ و بنابراین در $\Struct{M}$ برقرار است، $x = n$ و در
  نتیجه $x$ استاندارد می‌بود. حالت
  $y = x \oplus x \oplus \mathbf{z}^*$ نیز مشابه است.
\item با استدلالی مشابه، بزرگ‌ترین بلوک وجود ندارد.
\item ترتیب بلوک‌ها چگال است: اگر $[x]$ از $[y]$ کوچک‌تر باشد، که در
  آن $x \not\approx y$، آنگاه بلوکی مانند $[z]$ میان آن دو وجود دارد
  که از هر دو متمایز است. فرض کنید $x \prec y$. مانند پیش،
  $x \oplus y$ بر دو بخش‌پذیر است، شاید با باقیمانده؛ پس یک
  $u \in \Domain{M}$ وجود دارد به‌طوری‌که یا
  $x \oplus y = u \oplus u$ یا
  $x \oplus y = u \oplus u \oplus \mathbf{z}^*$. عنصر $u$ میانگین
  $x$ و $y$ و بنابراین میان آن‌هاست. فرض کنید
  $x \oplus y = u \oplus u$؛ حالت دیگر مشابه است. اگر
  $u \approx x$، آنگاه برای $n$ استانداردی:
  \[
  x \oplus y = x \oplus n \oplus x \oplus n,
  \]
  پس $y = x \oplus n \oplus n$ و $x \approx y$، برخلاف فرض. نتیجه
  می‌گیریم $u \not\approx x$. استدلالی مشابه
  $u \not\approx y$ را می‌دهد.
\end{enumerate}

پس بلوک‌های نااستاندارد مانند اعداد گویا مرتب‌اند: یک ترتیب خطی
!!a{enumerable} بدون دو سر تشکیل می‌دهند. از این رو، برای هر دو مدل
نااستاندارد !!{enumerable} از حساب درست، یعنی $\Struct{M}_1$ و
$\Struct{M}_2$، تحویل‌های آن‌ها به زبانی که تنها شامل $<$ و $=$ است
یکریخت‌اند. در واقع، یک یکریختی $h$ را می‌توان چنین تعریف کرد: بخش‌های
استاندارد $\Struct{M}_1$ و $\Struct{M}_2$ با مدل استاندارد
$\Struct{N}$ و بنابراین با یکدیگر یکریخت‌اند. بلوک‌های بخش نااستاندارد
خود مانند اعداد گویا مرتب‌اند و ازاین‌رو بنا بر
\olref[dlo]{thm:cantorQ} یکریخت‌اند؛ یک یکریختی میان بلوک‌ها را می‌توان
درون هر بلوک گسترش داد: در هر جفت بلوک، عناصری دلخواه را متناظر کنیم و
سپس تصویر جانشین $x$ در $\Struct{M}_1$ را جانشین تصویر $x$ در
$\Struct{M}_2$ بگیریم. توجه کنید که از این نتیجه \emph{نمی‌شود} که
$\mathfrak{M}_1$ و $\mathfrak{M}_2$ در زبان کامل حساب یکریخت‌اند
(زیرا یکریختی همواره نسبت به یک امضاست): می‌توان جمع و ضرب را بر
$\Domain{M_1}$ و $\Domain{M_2}$ به شیوه‌های ناهم‌ریخت تعریف کرد.
(این نکته همچنین از قضیه‌ای مشهور از وُت نتیجه می‌شود که شمار مدل‌های
شمارای یک نظریهٔ کامل نمی‌تواند دقیقاً~$2$ باشد.)

\begin{prob}
نشان دهید که در یک مدل نااستاندارد حساب نمی‌تواند بزرگ‌ترین بلوک وجود
داشته باشد.
\end{prob}

\begin{prob}
بگذارید $\Lang{L}$ !!{language} مرتبهٔ اولی باشد که $<$ تنها
!!{predicate} آن (جز $\eq$) است، و بگذارید
$\Struct{N}=(\Nat,<)$. همهٔ زیرمجموعه‌های متناهی یا هم‌متناهی
$\Struct{N}$ بدون پارامتر تعریف‌پذیرند. نشان دهید این‌ها \emph{تنها}
زیرمجموعه‌های بدون پارامتر تعریف‌پذیر $\Struct{N}$ هستند.

(راهنمایی: نخست بگذارید $\Obj{prc}(x,y)$ اختصاری برای فرمول
$\Lang{L}$ زیر باشد: «$x$ پیشینِ بی‌واسطهٔ $y$ است»:
\[
x<y \land \lnot \lexists[z][(x<z \land z<y)].
\]
اکنون، با هر زیرمجموعهٔ بدون پارامتر تعریف‌پذیر $\Struct{N}$ فرمولی
$!A(x)$ در $\Lang{L}$ متناظر است. برای هر چنین $!A$، جملهٔ $!D$ را
در نظر بگیرید:
\[
\lexists[x][\lforall[y][\lforall[z][((x<y \land x<z \land
  \Obj{prc}(y,z) \land !A(y)) \lif !A(z))]]].
\]
نشان دهید $\Sat{N}{!D}$ اگر و تنها اگر زیرمجموعهٔ تعریف‌شده به‌وسیلهٔ
$!A$ متناهی یا هم‌متناهی باشد.

حال بگذارید $\Struct{M}$ مدلی نااستاندارد و مقدماتاً هم‌ارز با
$\Struct{N}$ باشد. اگر $a\in\Domain{M}$ نااستاندارد است، عناصری
$b,c\in\Domain{M}$ بزرگ‌تر از $a$ بگیرید، به‌طوری‌که $b$ پیشینِ
بی‌واسطهٔ $c$ باشد. آنگاه خودریختی $h$ از تحویل ترتیبی
$(\Domain{M},\prec)$ وجود دارد که $h(b)=c$ (چرا؟). بنابراین اگر $b$
$!A$ را ارضا کند، $c$ نیز آن را ارضا می‌کند (چرا؟). پس $!D$ در
$\Struct{M}$ درست است و در نتیجه در $\Struct{N}$ نیز درست است. اما
این نتیجه می‌دهد که زیرمجموعهٔ $\Struct{N}$ تعریف‌شده به‌وسیلهٔ $!A$
متناهی یا هم‌متناهی است.)
\end{prob}

\end{document}
